% =====================================================================
%  1.5 Visão geral da implementação
% =====================================================================
\section{Visão geral da implementação do Subcomponente 3.2.b}
\label{sec:arranjo}

A implementação se organiza em três blocos (\autoref{fig:fluxograma-geral}):
estruturação institucional e social, implantação física e entrega com
monitoramento. Duas premissas orientam a sequência. As comunidades são
definidas pelo trabalho técnico-social, e o modelo de gestão é pactuado antes
das obras \cite{mop2026}. O financiamento do BIRD é de R\$ 183,0 milhões
(cerca de US\$ 33,5 milhões), e a contrapartida da Sanepar, de R\$ 67,1 milhões
\cite{custo2026}.

\begin{figure}[htbp]
  \centering
  \caption{Visão geral da implementação do Subcomponente 3.2.b}
  \label{fig:fluxograma-geral}
  \fluxograma{fluxograma-geral}
  \fonte{Elaborado pelo IAT (2026) com base em \citeonline{mop2026}.}
\end{figure}

Toda contratação financiada pelo BIRD segue o ciclo da
\autoref{fig:ciclo-contratacao}. O IAT elabora o TR em cerca de três meses, o
envia ao Banco Mundial, conduz a licitação e assina o contrato. As atividades
de planejamento, seleção e articulação são feitas pela equipe interna do IAT.

\begin{figure}[htbp]
  \centering
  \caption{Ciclo das contratações financiadas pelo BIRD}
  \label{fig:ciclo-contratacao}
  \fluxograma{ciclo-contratacao}
  \fonte{Elaborado pelo IAT (2026) com base no POA 2026--2027.}
\end{figure}
